%=======================================================================
%  CHAPTER 5 — NUMERICAL PROBLEMS AND SOLUTIONS
%=======================================================================
\section*{\color{acc}Ch 5 \textperiodcentered{} Equalization and Diversity \gap
\textcolor{sub}{Numerical Problems \& Solutions}}
\vspace{-4pt}{\color{acc}\rule{\textwidth}{1.1pt}}\vspace{4pt}

{\color{sub}\footnotesize \mk{No PYQ in 22 papers has set a pure calculation on this chapter.}
What \emph{is} examined, repeatedly, is the \textbf{MSE optimum-weight derivation} \tF{3}, which
the papers ask for as bookwork with a diagram. The rest of this companion works the tutorial-set
problems, because they are the standard follow-ups and they make the diversity theory concrete.}

\lead{The formula sheet}
\begin{itemize}
  \item \textbf{MSE cost function} \;
        $\xi = E[x_k^2] + \boldsymbol{\omega}^T\mathbf{R}\boldsymbol{\omega}
             - 2\mathbf{p}^T\boldsymbol{\omega}$
  \item \textbf{Optimum weights} \;
        $\hat{\boldsymbol{\omega}} = \mathbf{R}^{-1}\mathbf{p}$, \qquad
        $\xi_{\min} = E[x_k^2] - \mathbf{p}^T\hat{\boldsymbol{\omega}}$
  \item \textbf{LMS update} \;
        $\boldsymbol{\omega}_{k+1} = \boldsymbol{\omega}_k + \mu\,e_k\,\mathbf{y}_k^{*}$
  \item \textbf{Selection diversity}, $M$ Rayleigh branches, mean SNR $\Gamma$ each:
        \[ \Pr[\text{SNR} \le \gamma] = \left(1 - e^{-\gamma/\Gamma}\right)^{M},
           \qquad \frac{\bar{\gamma}_M}{\Gamma} = \sum_{k=1}^{M}\frac{1}{k} \]
  \item \textbf{Maximal ratio combining} \;
        $\displaystyle \gamma_M = \sum_{i=1}^{M}\gamma_i$, so
        $\bar{\gamma}_M = M\,\Gamma$
  \item \textbf{Interleaver} \; degree $m$, word length $n$: an $(n,k)$ code handling bursts of
        $b < (n-k)/2$ becomes an $(mn, mk)$ code handling bursts of $mb$
\end{itemize}

\hr

\T{Problem 1 \quad Optimum tap weights for a 2-tap equalizer}
{\color{sub}\footnotesize \tF{3} \yr{\texttt{79 Bh} \m{8}, \textbf{79 Ch} \m{7},
\textbf{80 Ch} \m{7}} \gap$\cdot$\gap the papers ask for the \textbf{derivation}; this worked
example shows the derivation actually being used}

\begin{asked}{2079 Bhadra \textperiodcentered{} Q5 \textperiodcentered{} [8]
\ \textperiodcentered{} \ 2079 Chaitra \textperiodcentered{} Q6 \textperiodcentered{} [7]
\ \textperiodcentered{} \ 2080 Chaitra \textperiodcentered{} Q6 \textperiodcentered{} [7]}
Determine the optimal solution of the weights using MSE algorithm with its appropriate diagram
and derivation.
\end{asked}
\begin{asked}[PROBLEM]{not a PYQ \textperiodcentered{} numbers chosen to exercise the result \added}
A 2-tap adaptive equalizer sees an input correlation matrix
$\mathbf{R} = \begin{bmatrix} 1 & 0.5 \\ 0.5 & 1 \end{bmatrix}$, a cross-correlation vector
$\mathbf{p} = \begin{bmatrix} 0.8 \\ 0.5 \end{bmatrix}$ and a transmitted signal power
$E[x_k^2] = 1$. Find the optimum weight vector $\mathbf{w}_{opt}$ and the minimum mean square
error $\xi_{min}$.
\end{asked}

\lead{Step 1: why $\hat{\boldsymbol{\omega}} = \mathbf{R}^{-1}\mathbf{p}$ (the derivation)}
\begin{itemize}
  \item Error signal, with $\mathbf{y}_k$ the input vector and $x_k$ the known training symbol:
        \[ e_k = x_k - \mathbf{y}_k^T\boldsymbol{\omega} \]
  \item Square it and take expectations:
        \[ \xi = E[e_k^2]
               = E[x_k^2] + \boldsymbol{\omega}^T \mathbf{R} \boldsymbol{\omega}
                 - 2\,\mathbf{p}^T\boldsymbol{\omega} \]
        with $\mathbf{R} = E[\mathbf{y}_k\mathbf{y}_k^T]$ and
        $\mathbf{p} = E[x_k\mathbf{y}_k]$.
  \item $\xi$ is \textbf{quadratic} in $\boldsymbol{\omega}$, so the surface is a bowl with
        \textbf{one} minimum. Set the gradient to zero:
        \[ \nabla\xi = 2\mathbf{R}\boldsymbol{\omega} - 2\mathbf{p} = 0
           \;\Longrightarrow\;
           \boxed{\hat{\boldsymbol{\omega}} = \mathbf{R}^{-1}\mathbf{p}} \]
\end{itemize}

\lead{Step 2: invert $\mathbf{R}$}
\[ \det \mathbf{R} = (1)(1) - (0.5)(0.5) = 0.75 \]
\[ \mathbf{R}^{-1} = \frac{1}{0.75}\begin{bmatrix} 1 & -0.5 \\ -0.5 & 1 \end{bmatrix}
   = \begin{bmatrix} 1.3333 & -0.6667 \\ -0.6667 & 1.3333 \end{bmatrix} \]
{\color{sub}\footnotesize $\mathbf{R}$ must be \textbf{non-singular} for the solution to exist.
Here $\det\mathbf{R} = 0.75 \ne 0$ \checkmark}

\lead{Step 3: the optimum weights}
\[ \hat{\boldsymbol{\omega}}
   = \begin{bmatrix} 1.3333 & -0.6667 \\ -0.6667 & 1.3333 \end{bmatrix}
     \begin{bmatrix} 0.8 \\ 0.5 \end{bmatrix} \]
\[ w_0 = 1.3333(0.8) - 0.6667(0.5) = 1.0667 - 0.3333 = 0.7333 \]
\[ w_1 = -0.6667(0.8) + 1.3333(0.5) = -0.5333 + 0.6667 = 0.1333 \]
\[ \boxed{\hat{\boldsymbol{\omega}} = \begin{bmatrix} 0.7333 \\ 0.1333 \end{bmatrix}} \]

\lead{Step 4: the minimum MSE}
\[ \xi_{\min} = E[x_k^2] - \mathbf{p}^T\hat{\boldsymbol{\omega}}
              = 1 - \left[(0.8)(0.7333) + (0.5)(0.1333)\right] \]
\[ = 1 - (0.5867 + 0.0667) = 1 - 0.6533 = \boxed{0.3467} \]

\begin{itemize}
  \item \mk{Interpretation:} the equalizer has removed about 65\,\% of the error energy. The
        residual 0.3467 is what no linear 2-tap filter can remove.
  \item \textbf{Adding taps lowers $\xi_{\min}$}, but at the cost of processing time and slower
        convergence. That trade-off is the practical answer to ``how many taps?''.
  \item \mk{The examiner wants the derivation, not the arithmetic.} Reproduce Steps 1 and 2
        symbolically and show the bowl-shaped error surface; the numbers are the easy part.
\end{itemize}

\hr

\T{Problem 2 \quad Selection diversity: probability of a deep fade}
{\color{sub}\footnotesize \yr{tutorial set, and the standard companion to the
``why implement diversity'' question \tS{11}}}

\begin{asked}[PROBLEM]{not a PYQ \textperiodcentered{} the standard companion to the ``why diversity'' question \added}
Five branches of a selection diversity receiver each receive an independent Rayleigh fading
signal. The average SNR on every branch is 20 dB. Determine the probability that the SNR of the
selected branch drops below 10 dB, compute the mean SNR after selection, and compare both with a
single receiver having no diversity.
\end{asked}
{\color{sub}\footnotesize ``Why is there a need to implement diversity?'' is asked in 11 of the
22 papers; this is the arithmetic that answers it with a number. \added}

\lead{Step 1: form the ratio $\gamma/\Gamma$ \mk{(in linear terms, not dB)}}
\begin{itemize}
  \item Threshold $\gamma = 10\ \text{dB} = 10^{10/10} = 10$
  \item Mean per branch $\Gamma = 20\ \text{dB} = 10^{20/10} = 100$
  \item \[ \frac{\gamma}{\Gamma} = \frac{10}{100} = 0.1 \]
  \mk{Subtracting in dB gives the same answer, but watch the order:
  $10\ \text{dB} - 20\ \text{dB} = -10$ dB, and $10^{-10/10} = 0.1$. \checkmark
  Either route is fine; mixing them is not.}
\end{itemize}

\lead{Step 2: probability the SNR falls below the threshold}
\[ \Pr[\text{SNR} \le \gamma] = \left(1 - e^{-\gamma/\Gamma}\right)^{M}
   = \left(1 - e^{-0.1}\right)^{M} = (0.09516)^{M} \]

\penalty0\vspace{2pt}
\begin{center}\small
\begin{tabular}{@{}rll@{}}
\toprule
\hrow \thead{Branches $M$} & \thead{$\Pr[\text{SNR} \le 10\ \text{dB}]$} & \thead{Meaning} \\
\midrule
1 (no diversity) & $9.516\times10^{-2}$ & the link is below threshold \mk{9.5\,\% of the time} \\
2 & $9.056\times10^{-3}$ & about 1 in 110 \\
4 & $8.201\times10^{-5}$ & about 1 in 12\,000 \\
5 & $\boxed{7.804\times10^{-6}}$ & about 1 in 128\,000 \\
\bottomrule
\end{tabular}
\end{center}

\begin{itemize}
  \item \mk{Five-branch selection diversity improves the outage probability by more than four
        orders of magnitude} over a single receiver, from $9.5\times10^{-2}$ to
        $7.8\times10^{-6}$.
  \item That is the entire justification for diversity, expressed as a number.
\end{itemize}

\lead{Step 3: mean SNR after selection combining}
\[ \frac{\bar{\gamma}_M}{\Gamma} = \sum_{k=1}^{M}\frac{1}{k}
   = 1 + \frac{1}{2} + \frac{1}{3} + \frac{1}{4} + \frac{1}{5} = 2.2833 \]
\[ \bar{\gamma}_5 = 100 \times 2.2833 = 228.33 \;=\; \boxed{23.59\ \text{dB}} \]

\begin{itemize}
  \item So the mean SNR rises only \textbf{3.59 dB}, from 20 dB to 23.59 dB.
  \item \mk{The important point:} the gain in \textbf{mean} SNR is modest, but the gain in
        \textbf{outage probability} is enormous. Diversity works by cutting the
        \textbf{tail} of the distribution, not by lifting the average.
        \textbf{Say this in the theory answer too.}
  \item \textbf{Diminishing returns:}
        $M{=}2$ gives $1.5\Gamma$ (1.76 dB),
        $M{=}3$ gives $1.833\Gamma$ (2.63 dB),
        $M{=}4$ gives $2.083\Gamma$ (3.19 dB),
        $M{=}5$ gives $2.283\Gamma$ (3.59 dB).
        \mk{Most of the benefit is already in the second branch}, which is why practical systems
        use 2, occasionally 4, and almost never 5.
\end{itemize}

\lead{Step 4: compare with maximal ratio combining}
\begin{itemize}
  \item MRC output SNR is the \textbf{sum} of the branch SNRs, so
        \[ \bar{\gamma}_M = M\Gamma = 5 \times 100 = 500 = \boxed{26.99\ \text{dB}} \]
  \item MRC gains a full \textbf{6.99 dB} against selection diversity's 3.59 dB, and its outage
        probability at 10 dB falls to $7.67\times10^{-8}$, another factor of 100 better.
  \item \mk{The price} is $M$ full receivers plus accurate estimation of each branch's amplitude
        and phase. Selection diversity needs only to compare SNRs.
\end{itemize}

\hr

\T{Problem 3 \quad Interleaver design against burst errors}
{\color{sub}\footnotesize \tF{3} \yr{\textbf{\texttt{80 Bh}}, \textbf{\texttt{81 Bh}},
\texttt{81 Ba}} \gap$\cdot$\gap the papers ask for the \textbf{principle}; this shows the
design sum behind it}

\begin{asked}{2080 Bhadra \textperiodcentered{} Q5 (b) \textperiodcentered{} [4]
\ \textperiodcentered{} \ 2081 Bhadra \textperiodcentered{} Q7 \textperiodcentered{} [4]
\ \textperiodcentered{} \ 2081 Baishakh \textperiodcentered{} Q5 \textperiodcentered{} [2+6]}
Why interleaving is needed in wireless communication? Explain its working in brief.
\ \textperiodcentered{}\ Elaborate the working principle behind Interleaving with an appropriate
diagram. \ \textperiodcentered{}\ Define Interleaving.
\end{asked}
\begin{asked}[PROBLEM]{not a PYQ \textperiodcentered{} the design sum behind the principle \added}
An interleaver has $m$ rows of $n$-bit words. Each word carries $k$ source bits and $(n-k)$
parity bits of a block code. Show how the interleaver and coder together break up a burst of
channel errors, and size the interleaver for a $(15,11)$ block code with interleaver degree
$m = 8$: give the block size, the source bits per block, and the longest burst the arrangement
can correct.
\end{asked}

\lead{Step 1: what the block code alone can do}
\begin{itemize}
  \item A block code with minimum distance $d'$ corrects $t$ errors when $d' \ge 2t+1$.
  \item For an $(n,k)$ code the parity budget is $n-k$, so it can handle a burst of
        \[ b < \frac{n-k}{2} \]
  \item For the $(15,11)$ code: \; $n-k = 4$, so $t = 2$ and a burst of up to
        \textbf{2 bits} per codeword is correctable.
\end{itemize}

\lead{Step 2: what the interleaver does}
\begin{itemize}
  \item Coded bits are written into the array \textbf{by column} and read out \textbf{by row}.
  \item Two bits that were adjacent in a codeword end up \textbf{$m$ bit periods apart} on the
        channel.
  \item So a channel burst of length $l = mb$ is spread across $m$ different codewords, giving
        each of them only $b$ errors, which each can correct.
\end{itemize}

\lead{Step 3: the resulting combined code}
\begin{itemize}
  \item An $(n,k)$ code plus a degree-$m$ interleaver behaves as an
        \[ (mn,\; mk)\ \text{block code} \]
  \item For $(15,11)$ with $m = 8$:
        \[ (8 \times 15,\; 8 \times 11) = \boxed{(120,\,88)\ \text{block code}} \]
  \item Burst-handling capability:
        \[ l = m \times b = 8 \times 2 = \boxed{16\ \text{consecutive bits}} \]
  \item \mk{The code rate is unchanged} at $88/120 = 11/15 = 0.733$. The interleaver costs no
        redundancy at all: it buys burst protection purely with \textbf{delay and memory}.
\end{itemize}

\lead{Step 4: the delay penalty}
\begin{itemize}
  \item The interleaver must fill all $mn = 120$ bits before the first row can be sent, and the
        de-interleaver must fill 120 bits before the first row can be decoded.
  \item At a coded rate of, say, 22.8 kbps (GSM full rate) that is
        \[ \frac{2 \times 120}{22\,800} = 10.5\ \text{ms of added delay} \]
  \item \mk{Human speech tolerates delays below about 40 ms}, which is the hard cap on how deeply
        a voice system may interleave. Data services, with no such limit, interleave far more
        deeply.
\end{itemize}

\hr

\T{Problem 4 \quad Choosing the equalizer algorithm from the channel}
{\color{sub}\footnotesize \yr{the reasoning behind ``explain any two adaptive equalization
algorithms'' (73 Ma, 74 Ma)}}

\begin{asked}{2073 Magh \textperiodcentered{} Q5 (a) \textperiodcentered{} [4]
\ \textperiodcentered{} \ 2074 Magh \textperiodcentered{} Q10 (a) \textperiodcentered{} [5]}
Explain briefly adaptive equalization algorithms (any two).
\ \textperiodcentered{}\ Write short note on Adaptive Equalization.
\end{asked}
\begin{asked}[PROBLEM]{not a PYQ \textperiodcentered{} the reasoning that picks the algorithm \added}
A GSM mobile travels at 70 km/h on a 900 MHz carrier through a channel whose rms delay spread is
$\sigma_\tau = 5\ \mu$s. The symbol rate is 270.833 ksym/s. Decide how many equalizer taps are
needed, how often the equalizer must retrain, and hence which adaptive algorithm (LMS, RLS or
Kalman) suits the link.
\end{asked}

\begin{itemize}
  \item \textbf{Step 1: how many taps?} The tap span must cover the \textbf{maximum expected delay
        spread}, in symbol periods.
        \[ T_s = \frac{1}{270\,833} = 3.692\ \mu\text{s} \]
        \[ \frac{\sigma_\tau}{T_s} = \frac{5}{3.692} = 1.35\ \text{symbols} \]
        Allowing for the maximum excess delay being several times $\sigma_\tau$, an equalizer
        spanning \textbf{about 5 symbol periods} is needed. \mk{GSM equalizers are specified to
        handle 4 to 5 symbols of delay spread, which matches.}
  \item \textbf{Step 2: how fast does the channel change?}
        \[ v = 70\ \text{km/h} = 19.44\ \text{m/s}, \qquad
           \lambda = 0.3333\ \text{m}, \qquad f_m = \frac{v}{\lambda} = 58.33\ \text{Hz} \]
        \[ T_c \approx \frac{0.423}{f_m} = \frac{0.423}{58.33} = 7.25\ \text{ms} \]
  \item \textbf{Step 3: retraining rate.} The equalizer must converge and be used well inside
        7.25 ms.
        \[ \text{minimum retrain rate} = \frac{1}{7.25\ \text{ms}} = 138\ \text{times per second} \]
        \mk{GSM sends a training midamble in every 4.615 ms frame}, ie 217 times per second,
        comfortably inside the requirement. \checkmark
  \item \textbf{Step 4: which algorithm?}
  \begin{itemize}
    \item The channel changes every few milliseconds, and the GSM training sequence is only
          \textbf{26 bits} long.
    \item \textbf{LMS} typically needs hundreds of iterations to converge, so 26 bits is not
          enough.
    \item \textbf{RLS} converges in roughly $2N$ iterations for an $N$-tap filter, ie about 10
          for a 5-tap equalizer. \mk{RLS, or an MLSE with a channel estimator, is the only viable
          choice.}
    \item That is exactly why GSM receivers use \textbf{MLSE with Viterbi}, not a simple LMS
          transversal equalizer.
  \end{itemize}
\end{itemize}

\hr

\T{Problem 5 \quad Cell-edge SNR gain from two-branch diversity}
{\color{sub}\footnotesize \yr{the arithmetic behind ``how diversity improves the quality of the
network'' (\textbf{76 Bh} \m{2})}}

\begin{asked}{2076 Bhadra \textperiodcentered{} Q5 \textperiodcentered{} [2+4]}
Explain how diversity improves the quality of network. Explain the operation of RAKE receiver
with appropriate block diagram.
\end{asked}
\begin{asked}[PROBLEM]{not a PYQ \textperiodcentered{} the number that answers ``how much'' \added}
On a Rayleigh channel the mean SNR at the cell edge is 15 dB and the receiver needs 8 dB to
demodulate. Compare the outage probability with a single antenna and with two-branch selection
diversity, and state the coverage gain this implies.
\end{asked}

\begin{itemize}
  \item \textbf{Step 1: the ratio.}
        \[ \gamma = 10^{8/10} = 6.31, \qquad \Gamma = 10^{15/10} = 31.62,
           \qquad \frac{\gamma}{\Gamma} = 0.1995 \]
  \item \textbf{Step 2: outage with one antenna.}
        \[ \Pr = 1 - e^{-0.1995} = 0.1809 = \boxed{18.1\ \%} \]
  \item \textbf{Step 3: outage with two branches.}
        \[ \Pr = (0.1809)^2 = 0.03273 = \boxed{3.27\ \%} \]
  \item \textbf{Step 4: what that is worth in dB.}
        To reach 3.27\,\% outage with a \emph{single} antenna the mean SNR would have to satisfy
        $1 - e^{-\gamma/\Gamma'} = 0.03273$, ie $\gamma/\Gamma' = 0.03327$, giving
        \[ \Gamma' = \frac{6.31}{0.03328} = 189.6 = 22.78\ \text{dB} \]
        \[ \text{diversity gain} = 22.78 - 15 = \boxed{7.78\ \text{dB}} \]
  \item \mk{Interpretation:} a second antenna is worth almost 8 dB of transmit power at this
        operating point, for no extra spectrum and no extra power.
        \textbf{That is why every base station has two receive antennas.}
  \item \textbf{Coverage implication.} With a path loss exponent $n = 3.5$, 7.78 dB of margin
        extends the cell radius by
        \[ 10^{7.78/(10 \times 3.5)} = 10^{0.2223} = 1.67\times \]
        ie \mk{about 2.8 times the area from the same site}.
\end{itemize}
